\documentclass{article}

\usepackage{amsmath,amssymb}
\usepackage{xurl}
\usepackage[hidelinks]{hyperref}
\setlength{\emergencystretch}{2em}

% Shared commands for notes in docs/paper-gaps/.

\DeclareUnicodeCharacter{2080}{\ifmmode{}_{0}\else\textsubscript{0}\fi}
\DeclareUnicodeCharacter{2081}{\ifmmode{}_{1}\else\textsubscript{1}\fi}
\DeclareUnicodeCharacter{2082}{\ifmmode{}_{2}\else\textsubscript{2}\fi}
\DeclareUnicodeCharacter{2099}{\ifmmode{}_{n}\else\textsubscript{n}\fi}

\newcommand{\Lean}{\textsc{Lean}}
\ExplSyntaxOn
\NewDocumentCommand{\leanid}{m}{
  \ifmmode
    \textsf{\detokenize{#1}}
  \else
    \group_begin:
    \urlstyle{sf}
    \tl_set:Nn \l_tmpa_tl {#1}
    \tl_replace_all:Nnn \l_tmpa_tl {\_} {_}
    \exp_args:NV \path \l_tmpa_tl
    \group_end:
  \fi
}
\ExplSyntaxOff
\providecommand{\tr}{\operatorname{tr}}
\newcommand{\tnleanIssueBase}{https://github.com/LionSR/TNLean/issues/}
\newcommand{\tnleanPrBase}{https://github.com/LionSR/TNLean/pull/}
\newcommand{\tnleanDocBase}{https://sirui-lu.com/TNLean/docs/}
\newcommand{\ghissue}[1]{\href{\tnleanIssueBase#1}{GitHub issue \##1}}
\newcommand{\ghpr}[1]{\href{\tnleanPrBase#1}{PR \##1}}
\newcommand{\doclink}[1]{\href{\tnleanDocBase#1.html}{\path{#1}}}

% Dirac notation and the identity operator, as in blueprint/src/macros/common.tex.
% \providecommand keeps note-local definitions authoritative.
\usepackage{dsfont}
\providecommand{\ket}[1]{|#1\rangle}
\providecommand{\bra}[1]{\langle #1|}
\providecommand{\braket}[2]{\langle #1 | #2 \rangle}
\providecommand{\ketbra}[2]{|#1\rangle\!\langle #2|}
\providecommand{\Id}{\mathds{1}}
\providecommand{\one}{\mathds{1}}
\providecommand{\C}{\mathbb{C}}
\providecommand{\MN}[1]{M_{#1}(\C)}
\providecommand{\MD}{M_D(\C)}

% External referencing for paper sources.  The build helper writes sanitized
% auxiliary files under build/paper-aux/ and excludes duplicated source labels.
\usepackage{xr}

% Import a generated paper auxiliary file with a caller-supplied label prefix.
% The candidate paths support builds from docs/paper-gaps/, the repository root,
% and build/paper-aux/ itself.
\newcommand{\importpaperaux}[2]{%
  \IfFileExists{../../build/paper-aux/#2.aux}{%
    \externaldocument[#1]{../../build/paper-aux/#2}%
  }{%
    \IfFileExists{build/paper-aux/#2.aux}{%
      \externaldocument[#1]{build/paper-aux/#2}%
    }{%
      \IfFileExists{#2.aux}{%
        \externaldocument[#1]{#2}%
      }{}%
    }%
  }%
}

% General external reference macros
\newcommand{\paperref}[2]{\ref{#1-#2}}
\newcommand{\papereqref}[2]{(\ref{#1-#2})}

% CPSV16 source (1606.00608 / MPDO-22-12-17-2.tex) external document and shortcuts
\importpaperaux{cpsv16-}{1606.00608/MPDO-22-12-17-2-tnlean}
\newcommand{\cpsvref}[1]{\paperref{cpsv16}{#1}}
\newcommand{\cpsveqref}[1]{\papereqref{cpsv16}{#1}}

% Verdict marker: \gapnote{<kind>}{<status>}. Kind is the policy
% classification (clarification, local-correction, scope-restriction,
% unfaithful, false-source, open-gap); status is open, wip (elimination
% underway), resolved, or historical. Machine-read by the site index;
% prints a verdict line.
\newcommand{\gapnote}[2]{\par\noindent\textbf{Verdict:} #1 (#2).\par}


\title{Adjacent regions obstruct the unrestricted CID forward implication in
Cirac--Perez-Garcia--Schuch--Verstraete 2017}
\author{}
\date{2026-07-25}

\begin{document}
\maketitle
\gapnote{false-source}{resolved}

\section{Source statement and the adjacent gap}

Definition~3.3 of Ref.~\cite{Cirac2016MPDO_arXiv} defines correlations
independent of the distance by quantifying over observables on any two
disjoint regions.\footnote{Local source:
\path{Papers/1606.00608/MPDO-22-12-17-2.tex:437--448}, label \path{ZCL}.}
Disjoint regions may be adjacent. If $n_2=n_1'+1$, no free site lies between
the observables. The allowed shifts satisfy
\begin{align}
    n_1'-n_2+1
        &<\Delta<N-n_2',
    \label{eq:cpsv16_pure_zcl_adjacent_gap_cid_scope_allowed_shift_interval}
\end{align}
which in the adjacent case means $\Delta\geq 1$ and leaves at least one free
site on each complementary arc.

For single-site observables at positions $n_1<n_2$, the source computes the
two-observable expectation by\footnote{Local source:
\path{Papers/1606.00608/MPDO-22-12-17-2.tex:490--496}, label \path{Corr}.}
\begin{align}
    \bra{V}O_1O_2\ket{V}
        &=\tr\!\left(\mathbb E_{O_1}
          \mathbb E^{\,n_2-n_1-1}
          \mathbb E_{O_2}
          \mathbb E^{\,N+n_1-n_2-1}\right).
    \label{eq:cpsv16_pure_zcl_adjacent_gap_cid_scope_single_site_trace_formula}
\end{align}
For block observables on $[n_1,n_1']$ and $[n_2,n_2']$, the transfer exponents
are
\begin{align}
    g_{\mathrm{mid}}
        &=n_2-n_1'-1,
    \label{eq:cpsv16_pure_zcl_adjacent_gap_cid_scope_middle_gap}\\
    g_{\mathrm{wrap}}
        &=N+n_1-n_2'-1.
    \label{eq:cpsv16_pure_zcl_adjacent_gap_cid_scope_wrapping_gap}
\end{align}

The forward claim at line~498 of Ref.~\cite{Cirac2016MPDO_arXiv} states that
$\mathbb E^2=\mathbb E$ implies CID. Its argument substitutes
$\mathbb E^n=\mathbb E$ into the trace formula. Idempotence justifies
\begin{align}
    \mathbb E^n
        &=\mathbb E,
        \qquad n\geq 1,
    \label{eq:cpsv16_pure_zcl_adjacent_gap_cid_scope_positive_power_idempotence}
\end{align}
but not the corresponding identity at exponent zero. For adjacent regions,
\begin{align}
    g_{\mathrm{mid}}
        &=0,
    \label{eq:cpsv16_pure_zcl_adjacent_gap_cid_scope_adjacent_zero_gap}\\
    \mathbb E^{g_{\mathrm{mid}}}
        &=\mathbb E^0
    \label{eq:cpsv16_pure_zcl_adjacent_gap_cid_scope_zero_power_transfer}\\
        &=\Id.
    \label{eq:cpsv16_pure_zcl_adjacent_gap_cid_scope_zero_power_identity}
\end{align}
Idempotence does not identify $\Id$ with $\mathbb E$. The printed argument
therefore covers only configurations with at least one free site between the
observables, whereas Definition~3.3 also includes adjacent regions.

\section{Counterexample: the Bell-pair chain}

Take the canonical normal-tensor renormalization fixed point of
\path{lem:charact-NT-pure-RFP} at line~1300 of
Ref.~\cite{Cirac2016MPDO_arXiv}, with trivial gauge $X=\one$, trivial isometry
$U=\one$, and maximally mixed spectrum
\begin{align}
    \Lambda
        &=\left(\tfrac12,\tfrac12\right).
    \label{eq:cpsv16_pure_zcl_adjacent_gap_cid_scope_maximally_mixed_spectrum}
\end{align}
The tensor is
\begin{align}
    A^{(\alpha,\beta)}_{\gamma,\delta}
        &=\delta_{\gamma,\alpha}
          \delta_{\delta,\beta}\sqrt{\tfrac12},
        \qquad
        (\alpha,\beta)\in\{0,1\}^2.
    \label{eq:cpsv16_pure_zcl_adjacent_gap_cid_scope_bell_pair_tensor}
\end{align}
The physical index consists of two spins at each node. The generated states
are the chains described by source equations
\path{eq:III_MPSRFP}--\path{eq:III_varphi} at lines~564--578 of
Ref.~\cite{Cirac2016MPDO_arXiv}. Each node $n$ carries spins $a_n,b_n$, and
\begin{align}
    \ket{\varphi}
        &=\frac{\ket{00}+\ket{11}}{\sqrt2}
    \label{eq:cpsv16_pure_zcl_adjacent_gap_cid_scope_bell_state}
\end{align}
is shared between $b_n$ and $a_{n+1}$.

The transfer map is
\begin{align}
    \mathbb E(X)
        &=\tr(X)\operatorname{diag}(\Lambda)
    \label{eq:cpsv16_pure_zcl_adjacent_gap_cid_scope_transfer_trace_form}\\
        &=\tfrac12\tr(X)\one.
    \label{eq:cpsv16_pure_zcl_adjacent_gap_cid_scope_transfer_depolarizing_form}
\end{align}
Since
\begin{align}
    \tr(\operatorname{diag}(\Lambda))
        &=1,
    \label{eq:cpsv16_pure_zcl_adjacent_gap_cid_scope_spectrum_normalization}
\end{align}
the transfer map is idempotent:
\begin{align}
    \mathbb E^2
        &=\mathbb E.
    \label{eq:cpsv16_pure_zcl_adjacent_gap_cid_scope_transfer_idempotence}
\end{align}
The tensor is a single normal tensor: it is already one-block injective and
left-canonical. No repeated copies or copy weights occur.

Let $Z$ be the Pauli matrix. Place $Z$ on spin $b_1$ and another $Z$ on spin
$a_2$. These observables occupy adjacent regions with zero gap. Their
expectation is
\begin{align}
    \tr\!\left(\mathbb E_{I\otimes Z}
      \mathbb E^0
      \mathbb E_{Z\otimes I}
      \mathbb E^2\right)
        &=\bra{\varphi}Z\otimes Z\ket{\varphi}
    \label{eq:cpsv16_pure_zcl_adjacent_gap_cid_scope_adjacent_bell_expectation}\\
        &=1.
    \label{eq:cpsv16_pure_zcl_adjacent_gap_cid_scope_adjacent_value}
\end{align}
After the allowed shift $\Delta=1$, one free site lies on each complementary
arc. The expectation factorizes:
\begin{align}
    \tr\!\left(\mathbb E_{I\otimes Z}
      \mathbb E^1
      \mathbb E_{Z\otimes I}
      \mathbb E^1\right)
        &=\bra{\varphi}Z\otimes\one\ket{\varphi}^{2}
    \label{eq:cpsv16_pure_zcl_adjacent_gap_cid_scope_shifted_factorization}\\
        &=0.
    \label{eq:cpsv16_pure_zcl_adjacent_gap_cid_scope_shifted_value}
\end{align}
The two values differ. Thus the tensor does not satisfy CID in the unrestricted
sense of Definition~3.3 of Ref.~\cite{Cirac2016MPDO_arXiv}, although its
transfer map is idempotent.

Idempotence, normality, the adjacent value one, and the shifted value zero all
have kernel-checked proofs.\footnote{Formal conjunction:
\leanid{MPSTensor.bellPairChain_isNormalTensor_isTransferIdempotent_and_not_isPhysicalCID}
in \doclink{TNLean/MPS/RFP/BellPairCIDObstruction}.}

\section{Consequence for the theorem surface}

The unrestricted forward implication at line~498 of
Ref.~\cite{Cirac2016MPDO_arXiv} is false as printed. The counterexample is a
single normal renormalization fixed point. Unlike the raw-weight
counterexample in
\path{cpsv16_pure_zcl_raw_weight_counterexample.tex}, which refutes the
converse direction of Theorem~3.8, labelled \path{TheoremZCLPure}, by using
copy weights, this counterexample has no repeated copies. Normalizing copy
weights therefore cannot repair the forward implication; the quantification
over adjacent regions must be restricted.

The faithful repair requires a positive gap: each complementary arc must
contain at least one free site. Under this restriction, transfer idempotence
implies physical CID. The formal theorem is
\leanid{MPSTensor.isPositiveGapPhysicalCID_of_isTransferIdempotent}. The
Bell-pair chain satisfies the repaired statement because its positive-gap
configurations factorize, as
(\ref{eq:cpsv16_pure_zcl_adjacent_gap_cid_scope_shifted_factorization})
demonstrates.

The source-faithful unrestricted predicate
\leanid{MPSTensor.IsPhysicalCID} remains in the formal development with a
docstring pointer to this note. The present counterexample closes the question
left open by the raw-weight note, whose converse-direction counterexample did
not decide the forward implication for adjacent physical regions.

The verdict is a source claim that is false as printed, with a proved repair.
The printed argument fails at
(\ref{eq:cpsv16_pure_zcl_adjacent_gap_cid_scope_zero_power_identity}), and the
positive-gap statement is the correct replacement target for \ghissue{4772}.

\bibliographystyle{alpha}
\bibliography{references}

\end{document}